\section{命令接口：\texttt{cmd.h}}

\codefile{src/include/kernel/cmd.h}
\begin{lstlisting}[style=cstyle]
typedef int (*cmd_fn_t)(int argc, char** argv);

typedef struct {
    const char* name;   // 命令名（如 "cls"）
    const char* help;   // 帮助文本
    cmd_fn_t fn;        // 处理函数
} cmd_t;
\end{lstlisting}

\begin{figure}[H]
  \centering
  % figures/ch06/fig02-cmd_t.tex — auto-extracted from chapter source
\begin{tikzpicture}[node distance=0cm]
  % cmd_t 结构
  \node[memcell, minimum width=4cm, fill=blue!10, minimum height=0.7cm] (name) {name (const char*)};
  \node[memcell, minimum width=4cm, fill=green!10, minimum height=0.7cm, above=0cm of name] (help) {help (const char*)};
  \node[memcell, minimum width=4cm, fill=orange!10, minimum height=0.7cm, above=0cm of help] (fn) {fn (cmd\_fn\_t)};

  \node[left=0.3cm of name, font=\small\bfseries] {cmd\_t};

  % Pointers
  \node[right=2cm of name, font=\ttfamily\small] (namestr) {"cls"};
  \draw[pointer] (name.east) -- (namestr.west);

  \node[right=2cm of fn, font=\ttfamily\small] (fnaddr) {cmd\_cls\_main()};
  \draw[pointer] (fn.east) -- (fnaddr.west);
\end{tikzpicture}

  \caption{\texttt{cmd\_t} 结构体的内存布局与指针指向}
\end{figure}

每个命令做成独立编译单元（一个 \texttt{.c} 文件），在文件末尾导出一个
\texttt{const cmd\_t}：

\codefile{src/kernel/cmds/cmd\_cls.c}
\begin{lstlisting}[style=cstyle]
// cmd_cls.c
static int cmd_cls_main(int argc, char** argv) {
    (void)argc; (void)argv;
    printk("\x1b[2J\x1b[H");  // ANSI 清屏 + 光标归位
    return 0;
}

const cmd_t cmd_cls = {
    .name = "cls",
    .help = "Clear screen (serial ANSI)",
    .fn   = cmd_cls_main,
};
\end{lstlisting}

\begin{designbox}
为什么用 \texttt{const cmd\_t}（全局常量）而非注册函数？

\begin{enumerate}
  \item 编译期确定：不需要运行时动态分配内存（此时堆还没初始化）
  \item 链接器自动包含：只要 \texttt{.o} 文件被链接进来，\texttt{cmd\_t} 就存在
  \item Shell 只需 \texttt{extern} 声明 + 加入数组即可
\end{enumerate}
\end{designbox}

% ============================================================================

